\centerline{\bf \Huge Матпраздничный разнобой}

\begin{flushright}
        \scriptsize{}
\end{flushright}

\problem{1}
Расставьте в клетки квадрата 3×3 различные целые положительные числа, не большие 25, так, чтобы в любой паре соседних по стороне клеток одно число делилось на другое.

\problem{2}
Аня называет дату красивой, если все 6 цифр её записи различны. Например, 19.04.23 — красивая дата, а 19.02.23 и 01.06.23 — нет. А сколько всего красивых дат в 2023 году?

\problem{3}
Доктор Айболит хочет навестить и корову, и волчицу, и жучка, и червячка. Все четверо живут вдоль одной прямой дороги. Орлы готовы утром доставить Айболита к первому пациенту, а вечером забрать от последнего, но три промежуточных перехода ему придётся сделать пешком. Если Айболит начнёт с коровы, то длина его кратчайшего маршрута составит 6 км, если с волчицы — 7 км, а если с жучка — 8 км. Нарисуйте, как могли располагаться домики коровы, волчицы, жучка и червячка (достаточно одного примера расположения).

\problem{4}
Пять друзей подошли к реке и обнаружили на берегу лодку, в которой могут поместиться все пятеро. Они решили покататься на лодке. Каждый раз с одного берега на другой переправляется компания из одного или нескольких человек. Друзья хотят организовать катание так, чтобы каждая возможная компания переправилась ровно один раз. Получится ли у них это сделать?

\problem{5}
На столе лежат 6 яблок (не обязательно одинакового веса). Таня разложила их по 3 на две чашки весов, и весы остались в равновесии. А Саша разложил те же яблоки по-другому: 2 яблока на одну чашку и 4 на другую, и весы опять остались в равновесии. Докажите, что можно положить на одну чашку весов одно яблоко, а на другую два так, что весы останутся в равновесии.

\problem{6}
Одуванчик утром распускается, три дня цветет желтым, на четвертый день утром становится белым, а к вечеру пятого дня облетает. В понедельник днем на поляне было 20 желтых и 14 белых одуванчиков, а в среду — 15 желтых и 11 белых. Сколько белых одуванчиков будет на поляне в субботу?